\section{An exact dyadic determinant formula}\label{sec:valuation}
Individual residues in \eqref{eq:coeffmatrices} can have large dyadic
denominators. Adding their valuations cannot determine the denominator
of their sum. We instead express the whole coefficient matrix through
one finite orthogonal system. This identifies the cancellation exactly.

\subsection{Replacing residues by a finite binomial average}
Let $N=2K-1$ and let $L$ be the sum of $N$ independent variables taking
the values $1$ and $-1$ with equal weights. Expectations below are finite
rational sums. At an odd point $l=2j+1$,
\[
 O'(-l^2)=\frac{(-1)^j4^{K-1}(K-1-j)!(K+j)!}{l},\qquad
 \Pr(L=l)=\frac{N!}{2^N(K-1-j)!(K+j)!}.
\]
Pairing $l$ and $-l$ gives
\begin{equation}\label{eq:binomial}
 (B_K)_{mn}=\frac1{2N!}\E\left[
 \frac{L^4J(-L^2)}{(a_m-L^2)(a_n-L^2)}\right].
\end{equation}
Every denominator inside the expectation is odd and nonzero. Off the
diagonal the two denominator factors cancel factors of $J$; on the
diagonal a rational remainder survives. The argument below handles
both together.

\subsection{The finite orthogonal basis}
Define the monic Krawtchouk polynomials by
\begin{equation}\label{eq:krawgen}
 \sum_{n\geq0}p_n(L)\frac{z^n}{n!}
 =(1+z)^{(N+L)/2}(1-z)^{(N-L)/2}.
\end{equation}
At each support point, the right side is the product of $N$ factors
$1+z\varepsilon_i$. The expectation of the product of two such
generating functions is $(1+zs)^N$. Comparing coefficients and then
differentiating \eqref{eq:krawgen} gives
\begin{align}
 \E[p_np_r]&=\begin{cases}n!N!/(N-n)!,&n=r\leq N,\\0,&n\ne r,
 \end{cases}\label{eq:krawnorm}\\
 Lp_n&=p_{n+1}+n(N-n+1)p_{n-1},\qquad p_0=1,\ p_1=L.
 \label{eq:krawrec}
\end{align}
In particular $p_n$ is integral and has the parity of $n$.
Put $t=L^2$ and $q_j(t)=p_{2j}(L)$. The $K$ monic polynomials
$q_0,\ldots,q_{K-1}$ are orthogonal, with norms
\begin{equation}\label{eq:h}
 h_j=(2j)!N!/(N-2j)!.
\end{equation}
There are exactly $K$ distinct support points for $t$. The polynomial
$q_K=p_{N+1}$ vanishes at all of them, since the evaluated generating
function has degree $N$. Multiplication by $t$ is therefore represented
exactly by a finite $K$-by-$K$ matrix $T$ in this basis.

Applying \eqref{eq:krawrec} twice gives
\begin{align}
 tq_j&=q_{j+1}+\tau_jq_j+u_jq_{j-1},\label{eq:Trec}\\
 \tau_j&=(4j+1)N-8j^2,\qquad
 u_j=(2j)(2j-1)(N-2j+1)(N-2j+2).\nonumber
\end{align}
Terms outside the basis are omitted on the support. Every $\tau_j$ is
odd and every $u_j$ is divisible by four. Thus $T$ is triangular modulo
two with diagonal one. This supplies more than invertibility: it will
make every leading principal block of the relevant multiplier a unit.

\subsection{Valuation and nonvanishing for every subset}
For nonzero rational $r$, let $v_2(r)$ denote its exponent of two.
Write $\Z_{(2)}$ for the rationals with odd denominator. An odd
determinant is a unit in this ring.

\begin{theorem}\label{thm:valuation}
For every nonempty $S\subseteq\{1,\ldots,K\}$, with $d=|S|$, set
$\Delta_S=\prod_{m<n,\,m,n\in S}(4n^2-4m^2)$. Then
$\det(B_K)_S\ne0$ and
\begin{equation}\label{eq:valuation}
 v_2\bigl(\det(B_K)_S\bigr)
 =2v_2(\Delta_S)-d+
 \sum_{j=0}^{d-1}\left(v_2((2j)!)-v_2((2K-1-2j)!)\right).
\end{equation}
\end{theorem}
\begin{proof}
The purpose of the first change of basis is to replace the rational
functions in \eqref{eq:binomial} by polynomials of degree below $d$.
Put
\[
 W_S(t)=\prod_{m\in S}(a_m-t),\qquad
 F_S(t)=\frac{t^2J(-t)}{W_S(t)^2},\qquad
 r_m(t)=\frac{W_S(t)}{a_m-t}.
\]
Then $(B_K)_S$ is $(2N!)^{-1}$ times the Gram matrix
$\E[F_Sr_mr_n]$. The coefficient determinant of the $r_m$ in the
monomial basis is $\pm\Delta_S$: evaluate at the $a_m$ to obtain a
diagonal matrix and divide by the evaluation Vandermonde. Passing to
the monic basis $q_0,\ldots,q_{d-1}$ changes no further determinant.

Every $a_mI-T$ is invertible over $\Z_{(2)}$. Consequently
\[
 F_S(T)=T^2J(-T)W_S(T)^{-2}
\]
is defined over that ring and reduces to $T^{K+2-2d}$ modulo two.
Negative exponents are allowed because $T$ is invertible. This
reduction is triangular with diagonal one, so its leading $d$-by-$d$
block has odd determinant.

The functional calculus is finite: multiplication by $t$ acts on the
$K$ support points and $W_S$ vanishes at none of them. Orthogonality
therefore gives
\[
 \bigl(\E[q_iF_Sq_j]\bigr)_{i,j<d}
 =\diag(h_0,\ldots,h_{d-1})[F_S(T)]_{i,j<d}.
\]
The last determinant is a unit, hence nonzero. Including the basis
factor and the scalar $(2N!)^{-d}$ gives
$2v_2(\Delta_S)-d(1+v_2(N!))+\sum_{j<d}v_2(h_j)$.
Substitution of \eqref{eq:h} proves the formula.
\end{proof}

For example, $K=2$ and $S=\{2\}$ give
$B_{22}=-167/140$ and $C_{22}=1/5$. The theorem gives
$v_2(B_{22})=-2$. After making the pencil monic in $Z$, its $X$
coefficient is $-167/28$, again of valuation $-2$. The example
illustrates why the coefficient normalization must accompany the
matrix calculation.

\subsection{The denominator forced by the positive restriction}
From now on let
\begin{equation}\label{eq:selected}
 S_K=\{2,4,\ldots,2d\},\qquad d=\lfloor K/2\rfloor,\quad K\geq2,
\end{equation}
and define the polynomial monic in $Z$ and its integer clearer by
\begin{equation}\label{eq:monic}
 P_K(X,Z)=\det\bigl(C_S^{-1}(R_S+XB_S)+ZI_d\bigr),\qquad
 A_K=\min\{A\in\Z_{>0}:AP_K\in\Z[X,Z]\}.
\end{equation}
Since $P_K$ is monic, $A_KP_K$ is primitive. Its nonzero coefficient
of $X^d$ is $\det(C_S^{-1}B_S)$, so the reduced denominator of this
one coefficient must divide $A_K$.

Write $s_2(n)$ for the binary digit sum. The elementary factorial
formula $v_2(n!)=n-s_2(n)$ gives
$v_2\binom{2n}{n}=s_2(n)$. The selected Vandermonde is
\[
 \Delta_S=2^{2d(d-1)}\prod_{j=1}^d\frac{(2j-1)!}{j},
 \qquad v_2(\Delta_S)=3d^2+O(d\log(d+1)).
\]
Theorem~\ref{thm:valuation} and \eqref{eq:cpositive} now yield
\begin{align*}
 v_2(\det B_S)&=8d^2-2Kd+O(K\log K),\\
 v_2(C_{mm})&=4K-8+v_2(m)-s_2(K-m)-s_2(K+m),\\
 v_2(\det(C_S^{-1}B_S))&=8d^2-6Kd+O(K\log K)
 =-K^2+O(K\log K).
\end{align*}
We have proved the necessary bound
\begin{equation}\label{eq:necessaryA}
 \log A_K\geq(\log2)K^2-O(K\log K).
\end{equation}

The restriction matters. On the full coordinate space there is instead
the identity
\begin{equation}\label{eq:fullidentity}
 \det(C_K^{-1}B_K)=(-1)^K(2/7)^K\binom{2K}{K}^3.
\end{equation}
Indeed $J/O=1+\sum_l\rho_l/(u+l^2)$ implies, by divided differences
at $u=-a_m$ and the diagonal derivative,
$V\diag(\rho_l)V^T=-\diag(e_m)$. Taking determinants in
\eqref{eq:coeffmatrices} gives $(-1)^K\prod_l b_l/\prod_m c_m$,
which simplifies to \eqref{eq:fullidentity}. That full-space
coefficient has no dyadic denominator, but $C_K$ is indefinite.
Thus the necessary cost in \eqref{eq:necessaryA} belongs to the
particular positive restriction, not to every use of the functional.
